\section{Практическая работа №5}
\label{sec:pr5}

\subsection{Выполненные задания}

В пятой практической работе был создан Stack-проект \mintinline{text}{qc}. В проект добавлена библиотека QuickCheck и реализованы функции \mintinline{haskell}{addMod} и \mintinline{haskell}{reverseWords}. Для тестов создан каталог \mintinline{text}{test} и файл \mintinline{text}{Spec.hs}.

Перед свойствами были написаны генераторы:

\captionof{listing}{Генераторы для тестов}
\label{lst:generators-pr5}
\begin{minted}{haskell}
genWord :: Gen String
genWord = listOf1 $ elements ['a' .. 'z']

genSentence :: Gen String
genSentence =
  choose (0, 8) >>= \n ->
    replicateM n genWord >>= \wordsList ->
      return $ intercalate " " wordsList

genMod :: Gen Int
genMod = choose (1, 1000)
\end{minted}

Всего написано семь свойств: три для \mintinline{haskell}{addMod} и четыре для \mintinline{haskell}{reverseWords}. Тесты запускаются с подписями в консоли.

Основные коммиты практики: \mintinline{text}{7bf6419}, \mintinline{text}{ad439ae}, \mintinline{text}{928bdbe}, \mintinline{text}{462c309}, \mintinline{text}{c7d26e5}, \mintinline{text}{707ccf8}.

\subsection{Использование нейросетей}

История взаимодействия с LLM для пятой практики сохранена в файле \mintinline{text}{practice_05/qc/llm.txt}. Модель использовалась для обсуждения генераторов, свойств QuickCheck, типа \mintinline{haskell}{Arbitrary} и формы записи тестов.

\subsection{Вопросы защиты и их решения}

\mintinline{haskell}{Arbitrary} --- это класс типов QuickCheck для значений, которые можно автоматически генерировать. Если свойство имеет тип вроде \mintinline{haskell}{Int -> Int -> Bool}, QuickCheck берёт генераторы из представителей \mintinline{haskell}{Arbitrary} для \mintinline{haskell}{Int}. В работе также использовались собственные генераторы, чтобы ограничить модуль положительными значениями от 1 до 1000 и генерировать осмысленные строки из слов.

Функция \mintinline{haskell}{unwords} нужна для обратной сборки предложения из списка слов. \mintinline{haskell}{words} разбивает строку на слова, \mintinline{haskell}{reverse} меняет порядок слов, а \mintinline{haskell}{unwords} соединяет их обратно через пробел.

\newpage
